\section{Where architecture stops}
\label{sec:ceiling}

Table~\ref{tab:leaderboard} places the three systems that ran on the official
TabReD splits into the benchmark's own reference table, as if they were methods.
Table~\ref{tab:normalized} restates the same runs on a per-task normalised scale,
where $0$ is the weakest and $1$ the strongest of the eighteen published
baselines.

% нормированная таблица идёт первой намеренно: она маленькая и встаёт на
% ближайшую страницу, тогда как leaderboard занимает полосу целиком и
% вытесняет всё, что объявлено после него
\begin{table*}[t]
\centering
\caption{Normalised position on each TabReD task: $0$ is the weakest and $1$ the strongest of the 18 published baselines. The three agent rows share one backbone model, so the spread among them is attributable to architecture alone.}
\label{tab:normalized}
\small
\setlength{\tabcolsep}{4pt}
\begin{tabular}{lrrrrrrrrr}
\toprule
Method & Homesite & Ecom & HomeCred & Sberbank & Cooking & Delivery & Maps & Weather & Mean \\
\midrule
Terminus-2 & 0.88 & -0.03 & 0.93 & 0.77 & 0.60 & 0.97 & 0.83 & 0.80 & \textbf{0.72} \\
AutoDS-Tools & 0.95 & 0.48 & 0.97 & 0.72 & 0.82 & 0.98 & 0.92 & 0.95 & \textbf{0.85} \\
FEDOT.LLM & 0.91 & 1.43 & 0.92 & 0.96 & 0.22 & 0.69 & 0.46 & 0.66 & \textbf{0.78} \\
\midrule
XGBoost & 0.92 & 0.39 & 1.00 & 0.92 & 0.78 & 0.99 & 0.96 & 0.98 & 0.87 \\
MLP-PLR ens. & 1.00 & 0.91 & 0.93 & 1.00 & 1.00 & 0.89 & 1.00 & 0.89 & 0.95 \\
Linear & 0.00 & 0.16 & 0.59 & 0.73 & 0.03 & 0.77 & -0.00 & -0.00 & 0.28 \\
\bottomrule
\end{tabular}
\end{table*}

\begin{table*}[t]
\centering
\caption{Agent architectures placed in the TabReD leaderboard: full official temporal splits, the benchmark's metrics in raw units. Published rows are Optuna-tuned means over 15 seeds \citep{rubachev2024tabred}; agent rows use \texttt{gemma-4-31b-it}; agent rows are means over three attempts. Best per column in bold.}
\label{tab:leaderboard}
\small
\setlength{\tabcolsep}{3pt}
\begin{tabular}{lrrrrrrrrr}
\toprule
& \multicolumn{3}{c}{ROC-AUC $\uparrow$} & \multicolumn{5}{c}{RMSE $\downarrow$} & \\
\cmidrule(lr){2-4}\cmidrule(lr){5-9}
Method & Homesite & Ecom & HomeCred & Sberbank & Cooking & Delivery & Maps & Weather & Avg.\ rank \\
\midrule
\addlinespace[2pt]
\multicolumn{10}{l}{\itshape GBDT} \\
XGBoost & 0.9601 & 0.5763 & \textbf{0.8670} & 0.2419 & 0.4823 & 0.5468 & 0.1616 & 1.4671 & 5.4 \\
LightGBM & 0.9603 & 0.5758 & 0.8664 & 0.2468 & 0.4826 & 0.5468 & 0.1618 & \textbf{1.4625} & 6.6 \\
CatBoost & 0.9606 & 0.5596 & 0.8621 & 0.2482 & 0.4823 & \textbf{0.5465} & 0.1619 & 1.4688 & 7.0 \\
RandomForest & 0.9570 & 0.5764 & 0.8269 & 0.2640 & 0.4884 & 0.5959 & 0.1653 & 1.5838 & 17.8 \\
Linear & 0.9290 & 0.5665 & 0.8168 & 0.2509 & 0.4882 & 0.5579 & 0.1709 & 1.7679 & 19.0 \\
\addlinespace[2pt]
\multicolumn{10}{l}{\itshape Tabular DL} \\
MLP & 0.9500 & 0.6015 & 0.8545 & 0.2508 & 0.4820 & 0.5504 & 0.1622 & 1.5470 & 9.8 \\
SNN & 0.9492 & 0.5996 & 0.8551 & 0.2858 & 0.4838 & 0.5544 & 0.1651 & 1.5649 & 15.5 \\
DCNv2 & 0.9392 & 0.5955 & 0.8466 & 0.2770 & 0.4842 & 0.5532 & 0.1672 & 1.5782 & 17.0 \\
ResNet & 0.9469 & 0.5998 & 0.8493 & 0.2743 & 0.4825 & 0.5527 & 0.1625 & 1.5021 & 12.0 \\
FT-Transformer & 0.9622 & 0.5775 & 0.8571 & 0.2440 & 0.4820 & 0.5542 & 0.1625 & 1.5104 & 8.9 \\
MLP (PLR) & 0.9621 & 0.5957 & 0.8568 & 0.2438 & 0.4812 & 0.5527 & 0.1616 & 1.5177 & 6.9 \\
Trompt & 0.9588 & 0.5803 & 0.8355 & 0.2509 & 0.4809 & 0.5519 & 0.1624 & 1.5187 & 10.9 \\
\addlinespace[2pt]
\multicolumn{10}{l}{\itshape Ensembles} \\
MLP ens. & 0.9503 & 0.6019 & 0.8557 & 0.2447 & 0.4815 & 0.5494 & 0.1620 & 1.5186 & 7.8 \\
MLP-PLR ens. & \textbf{0.9629} & 0.5981 & 0.8585 & \textbf{0.2381} & \textbf{0.4806} & 0.5518 & \textbf{0.1612} & 1.4953 & 4.1 \\
\addlinespace[2pt]
\multicolumn{10}{l}{\itshape Training methodologies} \\
MLP aug. & 0.9523 & 0.6011 & 0.8449 & 0.2659 & 0.4832 & 0.5532 & 0.1631 & 1.5193 & 13.5 \\
MLP aug. rec. & 0.9531 & 0.5960 & 0.7453 & 0.2515 & 0.4834 & 0.5541 & 0.1636 & 1.5160 & 13.9 \\
\addlinespace[2pt]
\multicolumn{10}{l}{\itshape Retrieval-augmented} \\
TabR & 0.9487 & 0.5943 & 0.8501 & 0.2820 & 0.4828 & 0.5514 & 0.1639 & 1.4666 & 12.9 \\
ModernNCA & 0.9514 & 0.5765 & 0.8531 & 0.2593 & 0.4825 & 0.5498 & 0.1625 & 1.5062 & 11.9 \\
\midrule
\multicolumn{10}{l}{\itshape Agentic AutoML systems (this work)} \\
Terminus-2 \emph{(open harness)} & 0.9588 & 0.5585 & 0.8590 & 0.2491 & 0.4837 & 0.5482 & 0.1628 & 1.5247 & 11.6 \\
AutoDS-Tools \emph{(multi-agent)} & 0.9611 & 0.5800 & 0.8633 & 0.2515 & 0.4820 & 0.5474 & 0.1620 & 1.4783 & 6.9 \\
FEDOT.LLM \emph{(rigid pipeline)} & 0.9597 & \textbf{0.6200} & 0.8569 & 0.2398 & 0.4867 & 0.5620 & 0.1664 & 1.5651 & 11.9 \\
\bottomrule
\end{tabular}
\end{table*}


\paragraph{Architecture is worth about half the range.}
With the backbone held fixed, the normalised mean of Table~\ref{tab:normalized}
moves from $0.26$ for the rigid pipeline to $0.85$ for the multi-agent system as
deployed, a spread of $0.59$,
with the open harness at $0.72$ between them; in average rank, $18.2$ against
$6.6$ among twenty-one methods, with the open harness at $11.4$. Part of that
spread is prescription rather than architecture: the multi-agent system without
its instruction layer scores $0.75$, so the span attributable to architecture
alone is $0.49$. We use the smaller figure whenever the claim is about
architecture and the larger only where the systems are compared as they ship.

\paragraph{But it does not reach competent practice.}
Tuned XGBoost sits at $0.87$ (rank $5.1$) and the strongest published method at
$0.95$ (rank $4.1$). The best of our systems reaches $0.85$ at rank $6.6$, which
places it inside the gradient-boosting band rather than below it: LightGBM is at
$0.85$ and CatBoost at $0.80$, so the multi-agent system is level with the first,
ahead of the second, and short of XGBoost by $0.02$. The honest statement is a
narrower boundary than the usual one: \emph{with an open 31B model, the best
architecture we tested matches the middle of competently tuned gradient boosting
on real-world tabular tasks under temporal shift, but not its best member, and no
architecture we tested approaches the strongest published method.} We state it
plainly because the comparison is usually made against defaults, where the same
systems would appear to win outright.

\paragraph{The ordering is not monotonic in either direction.}
The system that ranks highest is neither the most nor the least constrained.
Read as it stands the relationship is an inverted U---rigid pipeline $0.26$,
multi-agent $0.85$, open harness $0.72$---but each system enters in the
configuration it ships in, and those differ in more than architecture: the
multi-agent row carries its prescribed-knowledge layer and an image provisioned
with the libraries that layer names, while the open harness receives the task
statement alone in the stock image. Two arms separate those strands, and both
have now been run. Provisioning alone does nothing
(Section~\ref{sec:matched}); prescription does almost everything. With its layer
removed---the \emph{neither} column of Table~\ref{tab:grid}, same tasks, same
image---the multi-agent system scores $0.75$ against $0.85$ as deployed, with the
open harness at $0.72$. The ordering survives the separation but the distance
does not: $0.13$ as deployed becomes $0.03$ once neither system carries a
prescription, so on these eight tasks the layer is worth roughly three times the
architectural difference between the two. Neither the comparison with the rigid
pipeline at $0.26$ nor the gap to XGBoost is affected.

\subsection{Provisioning alone does nothing}
\label{sec:matched}

We ran the open harness on the image built for the multi-agent system, same
eight tasks, three attempts each. Seven of the eight land inside this system's
own noise floor---relative differences from $+0.08\%$ to $-0.93\%$, with four of
them at or below $0.08\%$ and several reproducing the stock-image attempts to
every digit printed. The reason is in the traces rather than the metrics: across all twenty-four
trials the agent named none of the three libraries the enriched image
pre-installs. It went to gradient boosting by hand, as on the stock image, where
the library was already reachable because it had installed it itself.
Provisioning without a text that points at it changes the container and not the
run.

That the enriched image was itself defective makes the point sharper. It
installed the prescribed libraries on a base without the OpenMP runtime they link
against, so LightGBM imported cleanly and then failed inside
\texttt{fit\_predict} as a broken joblib worker pool---a failure that says
nothing about OpenMP on its face. Forty-five trials across the two systems met
it; forty-five diagnosed it, reasoned from the dead worker pool to the missing
\texttt{libgomp.so.1}, installed it and continued. None abandoned the library and
none reported the repair. The defect is ours and is recorded as a threat
(Section~\ref{sec:threats}); the behaviour is what belongs here. Provisioning is
inert as a treatment not because these agents ignore the environment but because
they rebuild it, and a system that will install what it needs cannot be studied
by installing it in advance.

The one exception is \texttt{sberbank-housing}, at $-3.49\%$ with a spread of
$0.017$ against $0.000$ to $0.001$ elsewhere: two attempts reproduce the stock
values exactly and the third trained a thousand boosting rounds without early
stopping, which under temporal shift cost that attempt $12\%$ of its RMSE and the
three-attempt mean the $3.49\%$ above. That is
the agent's own lapse and the subject of Section~\ref{sec:mechanism}.

\paragraph{The rigid pipeline sits at the bottom of the leaderboard.}
FEDOT.LLM's normalised mean is $0.26$ and its average rank $18.2$ of $21$, below
plain linear regression at $0.28$ and rank $18.6$: the most constrained of the
three systems, driving a mature AutoML backend, does not recover the score of the
simplest reference in the table. Its metrics reached us rounded to two decimals,
so we checked how much of the placement is an artefact of that rounding: sweeping
every task independently to both ends of its interval moves the normalised mean
between $0.05$ and $0.47$ and the rank between $20.3$ and $16.6$. The ordering
against linear regression is therefore \emph{not} resolvable at the precision
available to us and we rest nothing on it. What survives the whole interval is
what we do use: even at its most favourable corner the rigid pipeline stays
$0.25$ of the normalised range below the open harness, $0.37$ below the
multi-agent system, and below every tuned gradient-boosting baseline. We report
this about one of our own systems deliberately---it is the lower anchor of the
architectural axis, and its position is what gives the axis its length. Full
precision would need the original job artefacts of a system we could not rebuild,
so we could neither re-run it nor re-score its submissions.

\begin{table}[!t]
\centering
\caption{Cost and wall-clock per trial from the Harbor telemetry: means over trials of the trials' own cost, tokens and time. MLAB abbreviates MLAgentBench. Bold: the lowest value in its column within a benchmark.}
\label{tab:cost}
\scriptsize
\setlength{\tabcolsep}{2.5pt}
\begin{tabular}{llrrrrr}
\toprule
System & Bench & $n$ & \$/tr. & In/tr. & Out/tr. & Min/tr. \\
\midrule
Terminus-2 & TabReD & 24 & 0.0069 & 49\,969 & 2\,111 & \textbf{5.1} \\
AutoDS & TabReD & 24 & 0.0074 & 61\,606 & 5\,349 & 44.4 \\
FEDOT.LLM & TabReD & 24 & \textbf{0.0003} & \textbf{2\,680} & \textbf{155} & 33.1 \\
AutoDS ($K$ off) & MLAB & 30 & \textbf{0.0089} & \textbf{76\,463} & \textbf{5\,864} & 63.2 \\
AutoDS (shipped $K$) & MLAB & 30 & 0.0103 & 90\,439 & 6\,407 & 70.5 \\
Terminus-2 & MLAB & 24 & 0.4317 & 1\,989\,443 & 11\,873 & \textbf{19.1} \\
\bottomrule
\end{tabular}
\begin{tablenotes}\footnotesize
\item Distributions are right-tailed: the untreated AutoDS-Tools branch on MLAgentBench has a median of $10.9$ minutes against a mean of $63.2$, and one Terminus-2 trial (\texttt{feedback}, $571$ steps, \$8.71) is $84\%$ of its job, whose median trial cost \$0.033. FEDOT.LLM calls the model from inside its container, so its cost is its token count at list price.
\end{tablenotes}
\end{table}


\paragraph{Cost, and what actually binds.}
Table~\ref{tab:cost} reports the Harbor telemetry behind these results.

Token spend is not the binding constraint at this scale: a whole TabReD trial
costs the open harness under a cent and the most expensive row is under two. What
separates the systems from the tuned baselines is training compute and search,
not the price of the backbone. Nor do the orderings of accuracy and cost agree,
and they disagree by different amounts in different columns. On the common
adapter the higher-ranking system costs $1.07\times$ the dollars of the open
harness and $1.23\times$ its input tokens, which is nothing, against
$2.5\times$ its output tokens and $8.7\times$ its wall-clock, which is not.
Accuracy is bought here with the clock rather than with the bill.

The price of the layer is better read off the crossed cells than off the two
earlier AutoDS-Tools rows. Those two differ only in the instruction file and put
the layer at $+48\%$ input tokens, $+26\%$ output and $+43\%$ dollars, but they
ran on the earlier preparation and on an image already carrying the libraries the
prescription names. On the common adapter the money goes the other way and only
the time agrees: Table~\ref{tab:grid} has the bill falling from $0.0086$ to
$0.0074$ dollars a trial while the median trial rises from $15.4$ minutes to
$38.9$. Naming a library shortens the search and lengthens the fit, so the sign
of its cost depends on which resource is counted
(Section~\ref{sec:mechanism}).

What binds is the clock. In our re-run, twenty-four trials under the benchmark's
own $3600$\,s ceiling, agent wall-clock had a median of $38.9$ minutes and a
minimum of $18.7$; four trials finished within fifteen minutes of the ceiling and
four more were cut by it. A killed trial and a lost trial are not the same thing:
three of the four had already written a submission and are scored like any other,
and one had not, so the ceiling truncated the work in one trial of six and
destroyed the result in one of twenty-four. The open harness on the same tasks
and ceiling has a median of $3.4$ minutes and a maximum of $9.3$. We report the
lost trial as a failure rather than a zero, since a timeout is a property of the
budget and not a measurement of the model (Section~\ref{sec:threats}), and note
that because truncation falls on the slowest tasks it is not random with respect
to the result. An architecture averaging forty-one minutes where another averages
four is buying its accuracy partly with a margin it does not have.

The hour binds rather than the hardware, which a repeat shows. We ran this arm
twice, on a twelve-core workstation at concurrency two and on an eighteen-core
server at concurrency four: medians $37.4$ and $38.9$ minutes, so doubling the
parallel load and changing the machine moved the distribution by a minute and a
half. What moved was the loss rate, two trials of twenty-four against one,
because a faster machine leaves more agents with something on disk when the clock
stops. A timeout rate is a property of a deployment, not of a system. The
wall-clock column should be read accordingly: for the earlier jobs it is
end-to-end time divided by trials and concurrency, so it absorbs environment
builds and queueing, and the jobs did not share a machine. It is reliable as an
order of magnitude and as a within-pair comparison, not as a benchmark of
throughput.
